% Source: https://artofproblemsolving.com/community/c3h2016936_minimum_absolute_value_question
% Stable identity: aops-topic-2016936-q01
By the triangle inequality $|u|+|v|\ge|u-v|$,
\[
|x-a|+|x-d|\ge|(x-a)-(x-d)|=d-a,
\]
with equality exactly when $x-a$ and $x-d$ have opposite signs or one is zero, i.e. when $a\le x\le d$. Similarly
\[
|x-b|+|x-c|\ge c-b,
\]
with equality exactly when $b\le x\le c$. Adding,
\[
|x-a|+|x-b|+|x-c|+|x-d|\ge (d-a)+(c-b),
\]
and equality holds for every $x$ with $b\le x\le c$ (then also $a\le x\le d$). So the minimum value is
\[
\boxed{c+d-a-b},
\]
attained exactly for $b\le x\le c$.

The posted answer, minimum $0$ at $x=a$, is wrong: at $x=a$ the sum equals $(b-a)+(c-a)+(d-a)$, which is positive, and a sum of distances to four distinct points can never be $0$.
